\chapter{힉스 보손의 이해}

\lectureoverview{힉스 메커니즘은 구성 요소들이 필요해지는 순서대로 살펴보면 가장 잘 이해할 수 있다. 양자화된 운동과 장에서 출발해 평범한 그릇 모양 퍼텐셜을 멕시코 모자 퍼텐셜로 바꾸고, 그 결과 생기는 응축을 통해 페르미온과 약한 보손이 질량을 가질 수 있는 이유를 이해한다. 그런 다음에야 물리적 힉스 보손을 그 응축의 반지름 방향 진동으로 규정하고, 2012년에 그 존재를 밝혀낸 생성 및 붕괴 과정을 따라간다.}

\section{헤드라인 뒤에 있는 메커니즘}

$125\,\mathrm{GeV}$ 부근에서 새로운 입자가 발견되었다는 소식은 당연히 큰 흥분을 불러일으켰지만, 흥분이 곧 이해를 뜻하지는 않는다. 유용한 질문은 힉스가 중요한지, 또는 우주의 기원과 수사적으로 연결할 수 있는지가 아니다. 훨씬 더 구체적인 질문은 이 현상의 작동 원리가 무엇이며, 표준 모형에 왜 이것이 필요했는가이다.

명칭에 대해서도 간단히 덧붙일 점이 있다. Peter Higgs는 이 물리를 밝혀낸 여러 사람 가운데 한 명이며, Fran\c{c}ois Englert와 Robert Brout도 핵심 발상에 독립적으로 도달했다. 따라서 \emph{브루트----엔글레르}힉스 메커니즘이라는 표현이 역사적 사실을 더 충실히 반영하지만, 아래에서는 더 짧은 명칭을 사용한다.

설명은 여러 부분으로 나뉜다. 먼저 양자역학의 기본 개념 두 가지와 고전장론의 한 가지 개념이 필요하다. 다음으로 가장 낮은 에너지 상태에 영이 아닌 장이 들어 있는 진공이 필요하다. 그런 다음 표준 모형의 입자들을 살펴보고, 이론의 맨 형태에서는 왜 관련 입자들이 질량이 없는지 알아본다. 힉스 보손 자체는 거의 마지막에 등장한다. 힉스 보손이 전체 메커니즘인 것은 아니다. 메커니즘이 자리 잡은 뒤 남는 특정한 진동 가운데 하나다.

\section{양자역학의 두 가지 사실과 장}

첫 번째 양자역학적 사실은 양자화이다. 각운동량은 이산적인 단위로 나타나므로 좋은 전형이 된다.
\begin{equation}
L=n\hbar,
\qquad
n=0,\pm1,\pm2,\ldots .
\end{equation}
반정수 각운동량도 존재하지만, 지금의 논의에서는 이에 의존하지 않는다. 중요한 점은 연속적으로 변하는 고전적 운동이 양자 상태의 사다리로 바뀐다는 것이다.

두 번째 사실은 불확정성 원리이다. 가능한 진공들이 연속적인 원을 이루는 상황이 나타날 때까지 이 원리는 미뤄 두겠다. 그때가 되면 진공을 한 지점에 날카롭게 고정되어 정지해 있는 작은 고전적 공처럼 취급할 수 없게 된다.

이 두 양자역학적 사실 사이에는 고전적인 장 개념이 놓여 있다. 장은 공간과 시간의 모든 지점에 어떤 양을 부여한다. 전기장, 자기장, 중력장은 익숙한 예다. 겉보기에는 비어 있는 영역에도 장이 있을 수 있다. 서로 너무 멀리 떨어져 서로 보이지 않는 두 전하판을 떠올려 보자. 그 사이의 거의 균일한 전기장은 그래도 국소적으로 실재한다. 그러므로 ``빈 공간''이 ``모든 장이 0이라는 뜻일 필요는 없다.'' 이는 장들이 가능한 상태 가운데 가장 낮은 에너지 상태에 있다는 뜻이다.

장 에너지가 어떤 한 값에서 최소이고, 장이 그 값에서 벗어나면 에너지가 증가한다고 하자. 장이 교란되면 최솟값 주위에서 진동한다. 양자역학은 이 진동을 양자화하며, 그 각각의 양자가 입자다. 강의 전체에서 사용할 장론의 대응 관계는 다음과 같다.
\begin{equation}
\text{particle}
\quad\longleftrightarrow\quad
\text{quantum of a field vibration}.
\end{equation}

\section{장 공간 속의 그릇}

하나의 실수 장 성분 $\phi$에 대해 가장 단순한 에너지 지형은 그릇 모양이다.
\begin{equation}
V(\phi)=\frac{1}{2}m^{2}\phi^{2}.
\end{equation}
그릇의 바닥은 $\phi=0$에 있다. 바닥 주위의 작은 진동은 정해진 진동수를 가지며, 그 양자는 질량 $m$인 입자처럼 거동한다.

이제 단일 장을 두 성분 $(\phi_1,\phi_2)$으로 바꿔 보자. 에너지는 이 내부 평면에서 반지름에만 의존할 수 있다.
\begin{equation}
r^{2}=\phi_1^{2}+\phi_2^{2},
\qquad
V=V(r).
\end{equation}
이는 보통 공간의 움푹 팬 곳이 아니라 \emph{장 공간} 위의 그릇이다. 두 성분의 쌍 $(\phi_1,\phi_2)$을 회전시켜도 에너지는 변하지 않는다. 강의에서 사용한 모자는 이 대칭성을 눈에 보이게 해 준다. 축을 중심으로 모자를 돌려도 모양은 그대로다.

\begin{figure}[tbp]
\centering
\includegraphics[width=0.95\textwidth]{lecture_144_figure_02.png}
\caption{장 공간의 회전 대칭성을 직관적으로 보여 주기 위해 사용한 평범한 대칭 그릇과 실제 모자.}
\lectureframe{00:10:02}
\end{figure}

장이 움직이면 이 비유는 단순한 기하학을 넘어선다. 장을 중심에서 밀어내고 그릇 둘레를 따라 움직이게 해 보자. 이 순환 운동은 장 공간에서 각운동량을 가진다. 실험실에서 물체가 회전하는 각운동량이 아니라 내부 각운동량이다. 이를 양자화하면 보존 전하의 모형이 된다.
\begin{equation}
Q=nq_0,
\qquad
n\in\mathbb{Z}.
\end{equation}
이 방정식은 개략적인 표현이다. 핵심 물리적 발상은 전하를 띤 장의 들뜸을 내부 평면에서의 회전으로 그려 볼 수 있다는 점이다. 양전하와 음전하는 서로 반대 방향의 회전에 해당하고, 회전하지 않는 상태는 전하가 0인 경우에 해당한다.

\section{모자를 뒤집으면}

평범한 그릇에는 최저점이 하나뿐이다. 힉스 물리는 모자를 뒤집을 때 시작된다. 이제 중심은 불안정하고, 원형 골짜기 전체가 같은 최저 에너지를 갖는다. 이를 나타내는 표준적인 수학적 형태는 다음과 같다.
\begin{equation}
V(\phi_1,\phi_2)
=
\lambda
\left(
  \phi_1^{2}+\phi_2^{2}-v^{2}
\right)^{2},
\qquad \lambda>0,
\end{equation}
최솟값들은 다음 조건을 만족한다.
\begin{equation}
\phi_1^{2}+\phi_2^{2}=v^{2}.
\end{equation}
\editorialnote{명시적인 사차 퍼텐셜은 강의에서 그린 멕시코 모자 기하를 표준적인 해석적 형태로 완성한 것으로, 칠판의 수식을 그대로 옮긴 것은 아니다.}

\begin{figure}[tbp]
\centering
\includegraphics[width=0.95\textwidth]{lecture_144_figure_03.png}
\caption{뒤집힌 모자 퍼텐셜의 중앙 봉우리와 원형 골짜기. 전체 칠판 그림은 기하의 대부분이 빠진 이전 크롭 이미지를 대체한다.}
\lectureframe{00:13:10}
\end{figure}

진공은 이제 장이 0인 상태가 아니다. 진공은 골짜기 위 어딘가에 놓여야 한다.
\begin{equation}
(\phi_1,\phi_2)
=
v(\cos\theta,\sin\theta).
\end{equation}
방정식은 회전 대칭을 유지하지만, 특정한 진공을 택하면 하나의 각도 $\theta$가 정해진다. 이것이 자발적 대칭 깨짐이다. 법칙은 대칭이지만, 선택된 최저 에너지 상태는 대칭이 아니다.

이제 질적으로 다른 두 종류의 운동이 있다. 반径방향으로 안팎으로 움직이면 퍼텐셜을 거슬러 올라가므로 에너지가 든다. 골짜기 둘레를 따라 움직이면 $\theta$만 변하며, 고전적으로는 퍼텐셜 에너지가 들지 않는다. 따라서 전하와 비슷한 각운동량을 상당히 많이 지니더라도 추가 퍼텐셜 에너지는 거의 없을 수 있다.

\begin{figure}[tbp]
\centering
\includegraphics[width=0.95\textwidth]{lecture_144_figure_04.png}
\caption{골짜기에 표시된 한 점은 연속적인 원을 이루는 동등한 최솟값 가운데 선택된 특정한 진공을 나타낸다.}
\lectureframe{00:16:14}
\end{figure}

\section{불확정성이 진공을 응축으로 바꾸다}

고전적으로는 장을 가장자리의 한 점에 놓고 각속도를 0으로 설정하면 문제가 끝났다고 생각할 수 있다. 양자역학은 그렇게 날카롭게 결론 내리는 것을 허용하지 않는다. 위치와 운동량은 다음 관계를 만족한다.
\begin{equation}
\Delta x\,\Delta p\gtrsim\frac{\hbar}{2}.
\end{equation}
원형 골짜기에서 이에 대응하는 내용은 장의 각도를 국소화하면 그 켤레 내부 각운동량이 불확정해진다는 것이다. 강의의 개략적 표현으로는 다음과 같다.
\begin{equation}
\Delta\theta\,\Delta L_{\mathrm{int}}
\gtrsim
\frac{\hbar}{2}.
\end{equation}
\editorialnote{각도 연산자에는 기술적인 미묘함이 있으므로, 이 관계는 전역적으로 정확한 연산자 항등식이 아니라 강의에서 제시한 불확정성 논증을 국소적으로 준고전적으로 나타낸 것이다.}

따라서 골짜기의 한 점 부근에 집중된 진공은 내부 전하가 정해진 상태가 아니다. 여러 전하 고유상태가 폭넓게 중첩되어 있다. 강의에서는 이를 조작적으로 확인하는 방법을 제시한다. 전하를 한 단위 더하거나 한 단위 빼도 상태는 거의 변하지 않는다. 개략적으로는 다음처럼 나타낼 수 있다.
\begin{align}
\ket{\Omega}
&\sim
\sum_n c_n\ket{n},
\\
c_{n+1}
&\simeq
c_n
\quad\text{within the occupied range}.
\end{align}
그런 다음 $n$을 한 단위 옮겨도 이미 들어 있는 항들의 표지만 대부분 바뀐다. 이것이 진공이 \emph{응축}이라는 말의 뜻이다. 진공에는 관련 양자가 정해지지 않은 수만큼 들어 있어 하나를 더하거나 빼도 관측 가능한 거시적 상태는 달라지지 않는다.

실제 전약 진공은 보통의 전하가 응축된 상태가 아니다. 그런 상태라면 빈 공간의 전자기적 거동은 크게 달라질 것이다. 그러나 초전도체는 전하 응축의 구체적인 예를 제공한다. 그러므로 초전도체는 추상적인 멕시코 모자 그림과 진공이 양자적으로 결맞은 구조를 가질 수 있다는 주장을 잇는 유용한 물리적 연결고리다.

이로써 첫 번째 단계를 마친다. 최저 에너지 상태가 장이 0인 곳에 있을 필요는 없으며, 양자 불확정성으로 인해 그 상태가 개별 양자를 더하거나 빼는 데 둔감해질 수 있음을 알았다. 이제 어떤 입자들에 이런 배경이 필요한지 살펴보자.

\section{맨 처음의 표준 모형}

기본 입자 질량의 자연스러운 상한 기준은 플랑크 질량이다. 국소화된 에너지를 충분히 높이면 중력을 무시할 수 없게 되고, 결국 그 대상은 보통의 기본 입자라기보다 블랙홀로 보는 편이 더 적절해진다. 수치로는 다음과 같다.
\begin{equation}
m_{\mathrm P}
=
\sqrt{\frac{\hbar c}{G}}
\approx
2.2\times10^{-8}\,\mathrm{kg},
\end{equation}
아주 작은 먼지 알갱이의 질량 정도다. 알려진 기본 입자들은 이보다 엄청나게 가볍다. 강의에서 사용한 규모 추산에 따르면, 이들의 질량은 플랑크 질량의 약 $10^{-17}$에 불과하다.

이 관찰만으로 더 무거운 기본 입자가 존재하지 않는다는 것이 증명되지는 않는다. 가속기는 주어진 에너지로 생성할 수 있을 만큼 가벼운 입자를 우선적으로 드러낸다. 따라서 알려진 입자 스펙트럼은 입자를 찾는 데 사용된 장치의 도달 범위에 따라 일부 선택된 것이다. 현재의 접근 범위를 넘어서는 더 무거운 입자가 많이 있을 수도 있다.

관측된 입자들은 페르미온과 보손으로 나뉜다. 페르미온에는 전자, 뮤온, 타우, 이들의 중성미자, 그리고 여섯 가지 쿼크 맛이 포함된다. 보손에는 광자, 글루온, 약한 $W$ 및 $Z$ 보손이 포함된다. 글루온은 쿼크를 결합해 양성자와 중성자를 만들므로 일반적인 원자핵의 구조에서 핵심적인 역할을 한다.

\begin{figure}[tbp]
\centering
\includegraphics[width=0.95\textwidth]{lecture_144_figure_05.png}
\caption{강의에서 제시한 표준 모형의 첫 입자 목록. 페르미온 및 보손 열 위에는 방대한 질량 범위가 그려져 있다.}
\lectureframe{00:24:22}
\end{figure}

이론은 단순한 목록이 아니라 허용되는 상호작용의 집합이다. 전자는 광자를 방출할 수 있다.
\begin{equation}
e\longrightarrow e+\gamma.
\end{equation}
쿼크는 전하와 색전하를 모두 가지므로 광자나 글루온을 방출할 수 있다.
\begin{equation}
q\longrightarrow q+\gamma,
\qquad
q\longrightarrow q+g.
\end{equation}
전자는 글루온을 방출할 수 없고, 중성미자는 광자도 글루온도 방출하지 않는다. 적절한 약한 양자수를 가진 페르미온은 $Z$을 통해 상호작용할 수 있다.
\begin{equation}
f\longrightarrow f+Z.
\end{equation}
$W$ 보손과 하전 전류 반응도 표준 모형의 일부지만, 강의에서는 논의를 이어 갈 수 있도록 중성 약한 보손 하나만 남기고 이들을 생략한다.

\begin{figure}[tbp]
\centering
\includegraphics[width=0.95\textwidth]{lecture_144_figure_06.png}
\caption{광자, 글루온, $Z$ 방출 꼭짓점 옆에 적힌 ``질량 없음''. 이 과정들은 응축을 도입하기 전의 맨 처음 이론을 요약한다.}
\lectureframe{00:29:20}
\end{figure}

광자는 질량이 없다. 이와 전반적으로 같은 종류의 대칭 깨지지 않은 게이지 이론에서는 글루온과 약한 게이지 보손도 질량이 없다. 수정되지 않은 전약 이론 역시 게이지 대칭을 유지하면서 페르미온 질량을 손으로 그냥 넣을 수는 없다. 한 가지 요소가 빠져 있다. 그것을 도입하기 전에, 더 단순한 설정에서 배경장이 마찰처럼 작용하지 않고도 입자 질량을 어떻게 바꿀 수 있는지 살펴보자.

\section{배경장이 질량을 바꾸는 방식}

물 분자를 전기 쌍극자로 생각해 보자. 한쪽 끝은 약한 양전하를 띠고 다른 쪽 끝은 약한 음전하를 띤다. 빈 공간에서는 서로 반대인 두 방향이 에너지 면에서 동등하다. 강의에서는 두 방향을 서로 다른 내부 상태로 구별해 따라갈 수 있도록 장난스럽게 ``물''과 ``위스키''라는 이름을 붙인다.

분자를 축전기 판 사이에 놓자. 쌍극자 모멘트가 $\bm d$이고 전기장이 $\bm E$라면 상호작용 에너지는 다음과 같다.
\begin{equation}
U=-\bm d\mathbin{\cdot}\bm E.
\end{equation}
정렬된 두 방향에 대해서는 다음과 같이 된다.
\begin{equation}
U_{\parallel}=-dE,
\qquad
U_{\mathrm{anti}}=+dE.
\end{equation}
아인슈타인의 관계식 $E=mc^2$은 에너지 차이를 질량 차이로 바꾼다.
\begin{equation}
m_{\mathrm{anti}}-m_{\parallel}
=
\frac{2dE}{c^{2}}.
\end{equation}
분자에 마찰을 일으킨 것은 아무것도 없다. 배경장이 존재하는 동안 한 내부 상태의 에너지는 다른 상태보다 클 뿐이며, 따라서 질량도 더 크다.

\begin{figure}[tbp]
\centering
\includegraphics[width=0.95\textwidth]{lecture_144_figure_07.png}
\caption{축전기 전기장 안에서 서로 반대 방향을 향한 분자들. 장은 항력을 만들지 않고 두 상태의 에너지를 갈라놓는다.}
\lectureframe{00:34:35}
\end{figure}

이 예를 통해 오해를 부르는 두 가지 비유를 정리할 수 있다. 힉스 장은 입자들이 헤치고 나아가야 하는 당밀이 아니며, 질량을 가진 입자는 물질이 앞에 쌓여 느려지는 제설차가 아니다. 공간적으로 균일한 장에서는 알짜 힘이 전혀 없을 수도 있다. 효과는 내부 에너지의 변화이지, 마찰도 운동량 손실도 아니다.

같은 예를 양자 언어로 다시 설명할 수도 있다. 고전적인 전기장은 결맞은 상태이며, 여기서는 광자의 응축으로 설명할 수 있다. 분자는 그 배경에 광자를 방출하거나 배경에서 광자를 흡수할 수 있다. 점유량이 큰 결맞은 상태는 양자 하나를 더하거나 빼도 거의 변하지 않으므로, 광자 선은 응축에서 끝날 수 있다. 칠판에 그린 십자 표시는 에너지 보존이 실패했다는 뜻이 아니라, 광자가 배경 속으로 사라진다는 뜻이다.

\begin{figure}[tbp]
\centering
\includegraphics[width=0.95\textwidth]{lecture_144_figure_08.png}
\caption{광자 응축으로 다시 그린 축전기 장. 쌍극자가 방출한 광자는 양자 하나가 더 들어와도 변하지 않는 배경에 흡수된다.}
\lectureframe{00:38:45}
\end{figure}

이것이 필요한 틀이다. 응축이 두 내부 상태 사이의 양자수 차이를 흡수할 수 있다면, 두 상태를 서로 섞을 수 있다. 그 결과 생기는 진동은 질량으로 나타날 수 있다.

\section{힉스가 모든 질량을 설명하지는 않는 이유}

이 틀을 적용하기 전에 한 가지 오해를 바로잡아야 한다. 질량을 가진 모든 물체가 힉스 장에서 질량을 전부 얻는다는 말은 틀렸다.

완전히 반사하는 상자에 복사를 가두자. 개별 광자는 질량이 없지만, 닫힌 상자는 정지 좌표계를 가지며 에너지를 저장한다. 상자의 질량은 다음만큼 증가한다.
\begin{equation}
\Delta M=\frac{E_{\mathrm{radiation}}}{c^{2}}.
\end{equation}
따라서 구성 입자들의 정지 질량이 없어도, 가둬진 운동 에너지와 장 에너지로 질량이 생길 수 있다.

양성자는 중요한 실제 사례다. 양성자는 강하게 상호작용하는 상태에 있는 쿼크, 반쿼크, 글루온으로 이루어져 있다. 양성자 질량의 대부분은 이들의 운동 에너지와 결합 에너지에서 나오며, 구성 쿼크 몇 개에 힉스장이 부여한 질량을 단순히 더한 값에서 나오지 않는다. 강의에서 강조한 추산에 따르면, 가벼운 쿼크들의 질량을 영으로 만들어도 양성자 질량은 대략 퍼센트 수준만 달라진다. 블랙홀은 이 논점을 한층 더 분명하게 보여 준다. 블랙홀은 블랙홀을 형성한 전체 에너지 때문에 질량을 가질 수 있으며, 그 에너지가 모두 힉스장을 통해 먼저 기본 입자의 정지 질량으로 바뀌었기 때문일 필요는 없다.

힉스 메커니즘의 역할은 더 구체적이다. 바탕이 되는 게이지 이론의 일관성을 유지하면서 기본 페르미온과 약한 게이지 보손에 질량을 허용한다.

\section{키랄성: 질량은 왼손형과 오른손형의 혼합--}

먼저 전자에 주목하자. 스핀은 고유한 각운동량이다. 움직이는 입자에서는 스핀 방향과 운동량을 서로 비교할 수 있다. 강의의 기하학적 표현에 따르면, 오른손잡이 전자와 왼손잡이 전자는 이 두 방향의 관계로 구별된다. 질량이 없는 입자에서는 논증에 필요한 단순한 방식으로 헬리시티와 카이랄성이 일치하지만, 질량이 있는 입자에서는 두 개념을 정확히 구별해야 한다.

질량이 있는 전자를 계속 앞질러 간다고 상상해 보자. 다른 관성계에서는 전자의 운동량 방향이 뒤집히지만 스핀은 그대로이므로, 운동 방향을 기준으로 한 겉보기 손잡이성은 뒤집힐 수 있다. 어떤 관찰자도 빛의 속도로만 움직이는 엄밀히 질량이 없는 입자를 앞지를 수 없다. 이로부터 강의에서는 직관적인 설명을 제시한다. 질량이 없는 전자는 하나의 카이랄 상태를 유지하지만, 질량이 있는 전자는 좌우 상태를 끊임없이 섞는다.

정확한 장론의 표현에서는 디랙 질량항이 두 카이랄 장을 결합한다.
\begin{equation}
\mathcal L_m
=
-m_e\bar e e
=
-m_e
\left(
  \bar e_L e_R
  +
  \bar e_R e_L
\right).
\end{equation}
\editorialnote{디랙 질량항은 강의의 ``뒤집힘 속도'' 그림을 표준 수학 형태로 완성한 것이다. 카이랄성 혼합은 질량에 비례하지만, 실제 시간당 뒤집힘 횟수는 상태와 상대론적 시간 지연에 따라 달라진다.}

$m_e=0$이면 좌우 장은 서로 독립적으로 전파된다. $m_e\ne0$이면 방정식들이 두 장을 서로 결합한다. 이 정확한 의미에서 질량은 좌--우 혼합의 세기를 나타낸다.

\section{맨 처음 질량항을 금지하는 보존 법칙}

그렇다면 혼합항을 처음부터 그냥 쓰면 왜 안 되는가? 약한 상호작용은 $e_L$과 $e_R$을 서로 다르게 취급한다. Susskind는 이 차이에 일부러 단순한 이름을 붙인다. 바로 \emph{질크}다. 강의에서 사용한 정규화에서는 다음과 같다.
\begin{equation}
z(e_L)=1,
\qquad
z(e_R)=0.
\end{equation}
두 카이랄성은 전하는 같지만 약한 양자수는 다르다. 직접적인 전이
\begin{equation}
e_L\longrightarrow e_R
\end{equation}
는 약한 상호작용이 운반하는 양을 보존하지 못한다. 대칭이 깨지지 않은 이론에서는 맨 질량항이 금지된다.

\editorialnote{``''질크$SU(2)_L\times U(1)_Y$는 교육을 위한 약칭이다. 표준 모형에서는 왼손 및 오른손 페르미온이 서로 다른 $1$ 표현에 속하므로, 직접적인 디랙 질량항은 게이지 불변이 아니다. 강의에서 부여한 $0$와 }는 전약 상호작용의 전체 전하표를 제시하지 않고도 이 차이를 나타낸다.

칠판에서는 이 장애물을 연쇄 과정으로 보여 준다. 다른 무언가가 질크의 차이를 가져가지 않는 한 오른손 전자는 왼손 전자가 될 수 없다. 빠진 장이 도입되고 나면 종마다 결합 세기가 달라지지만, 같은 장애물이 쿼크, 뮤온, 타우에도 적용된다.

\begin{figure}[tbp]
\centering
\includegraphics[width=0.95\textwidth]{lecture_144_figure_09.png}
\caption{칠판에 그린 카이랄성의 연쇄 과정. 녹색 십자 위치에 응축 양자를 넣으면 질크 표지가 오른손 및 왼손 전파 상태 사이에서 바뀔 수 있다.}
\lectureframe{00:50:15}
\end{figure}

\section{``지그스'' 응축과 페르미온 질량}

정확히 빠져 있는 약한 양자수를 가진 양자가 생기는 장을 도입하자. 물리적 힉스 입자와 혼동되지 않도록 강의에서는 관련 양자를 \emph{지그스}라고 부른다. 지그스 장이 멕시코 모자 퍼텐셜을 가지며 그 진공이 응축이라고 하자. 그러면 왼손 전자는 지그스 양자를 방출해 오른손 전자가 될 수 있다.
\begin{equation}
e_L\longrightarrow e_R+\text{Ziggs},
\end{equation}
또는 오른손 전자가 지그스 양자를 흡수해 왼손 전자가 될 수 있다.
\begin{equation}
e_R+\text{Ziggs}\longrightarrow e_L.
\end{equation}
응축은 거시적 상태를 바꾸지 않고 양자수 차이를 흡수하거나 공급한다. 빈 대칭 깨지지 않은 진공에서는 금지되었던 일이 대칭이 깨진 진공에서는 허용된다.

반복되는 방출과 흡수는 $e_L$과 $e_R$을 서로 섞는다. 따라서 전자는 질량을 가진 입자로 전파된다. 같은 메커니즘이 하전 렙톤과 쿼크에도 적용된다. 각 종은 응축과 고유한 결합을 가지므로 뒤집힘의 세기, 따라서 질량도 종마다 다를 수 있다.

현대적 표기에서 게이지 불변 상호작용은 유카와 결합이다. 힉스 장이 진공기댓값 $v$를 얻으면 다음 항이 생긴다.
\begin{equation}
\mathcal L_Y
\supset
-\frac{y_f}{\sqrt{2}}
(v+h)\bar f f,
\end{equation}
따라서
\begin{equation}
m_f=\frac{y_fv}{\sqrt{2}},
\qquad
g_{hff}=\frac{y_f}{\sqrt{2}}
=
\frac{m_f}{v}.
\end{equation}
\editorialnote{이 수식은 강의의 응축 결합 논증에 해당하는 표준 전약 표기를 제시한다. 또한 같은 결합이 페르미온 질량과 힉스 보손의 해당 페르미온 붕괴율을 함께 결정하는 이유도 명시적으로 보여 준다.}

강의에서는 이를 카이랄 대칭의 자발적 깨짐으로 설명한다. 이는 예시에서 다룬 좌--우 이야기를 포착한다. 더 정확히 말하면, 힉스 진공기댓값이 전약 게이지 대칭을 전자기 부분군으로 깨뜨리고, 그 뒤 유카와 상호작용이 페르미온 질량항이 된다.

\section{약한 보손이 질량을 얻는 방식}

응축은 약한 게이지 보손도 바꾼다. 지그스 들뜸은 $Z$ 보손을 방출할 수 있고, 반대로 $Z$은 적절한 응축 지지 들뜸으로 바뀌었다가 다시 원래 상태로 돌아갈 수 있다. 따라서 $Z$은 힉스 장의 각도 방향 부문과 혼합되는 진공을 지나 전파된다. 그 결과 질량이 있는 벡터 보손이 생기며, 이 보손은 질량이 없는 벡터 보손에는 없는 종방향 편광을 가진다.

\begin{figure}[tbp]
\centering
\includegraphics[width=0.95\textwidth]{lecture_144_figure_10.png}
\caption{완성된 약한 보손 메커니즘의 칠판 그림. 위쪽에는 카이랄성 연쇄 과정이 남아 있고, 아래쪽에서는 $Z$이 응축 지지 부문과 양자수를 주고받는다.}
\lectureframe{00:53:40}
\end{figure}

게이지 보손의 경우 이것이 브루트--엔글레르--힉스 현상이다. 전체 표준 모형에서는 $W$ 보손과 $Z$이 질량을 얻으며, 원래의 중성 게이지 장 가운데 한 조합은 질량이 없는 광자로 남는다. 진공이 보통의 전자기 게이지 대칭을 깼다면, 초전도체 내부에서 실제로 일어나는 것처럼 광자에 질량이 생겼을 것이다. 그러면 빈 공간의 장거리 전자기 거동은 완전히 달라졌을 것이다.

물리적 힉스 입자가 분리되어 확인되기 훨씬 전부터 약한 보손은 이미 이 메커니즘을 강하게 뒷받침하는 증거를 제공했다. 약한 보손의 질량과 상호작용은 전약 대칭이 필요한 방식으로 깨졌음을 보여 주었다. 그래서 강의에서는 느슨한 의미에서 ``지그스'' 부문이 이미 알려져 있었다고 말할 수 있다. 실험에서 빠져 있던 대상은 독립적인 반径방향 진동이었다.

\editorialnote{강의에서 약 1980년에 SLAC에서 발견되었다고 즉흥적으로 언급한 대목은 여러 사건을 압축한 것이다. 중성 약한 전류는 1973년에 관측되었고, SLAC의 패리티 위반 전자 산란은 1970년대 후반에 보고되었으며, $W$ 및 $Z$ 보손은 1983년에 CERN에서 직접 발견되었다. 핵심 물리적 내용은 전약 대칭 깨짐이 그보다 앞서 실험적으로 확인되었다는 점이다.}

\section{힉스 보손의 실제 정체}

멕시코 모자로 돌아가 보자. 골짜기 둘레를 따라 움직이면 각도는 변하지만 반지름은 변하지 않는다. 각도 방향 부문은 질량이 있는 약한 보손에 흡수되는 모드를 제공한다. 물리적 힉스 보손은 다른 운동, 즉 퍼텐셜의 벽을 따라 안팎으로 움직이는 진동이다.

장을 개략적으로 다음과 같이 쓰자.
\begin{equation}
\Phi(x)
=
\frac{v+h(x)}{\sqrt{2}}
\exp\!\left(
  \frac{i\theta(x)}{v}
\right).
\end{equation}
변수 $\theta$는 골짜기를 따라 움직이는 것을 나타낸다. 변수 $h$은 반지름을 바꾼다. 게이지 이론에서는 적절한 각도 자유도가 질량이 있는 약한 보손의 종방향 성분이 되는 반면, $h$은 전파하는 스칼라 입자인 힉스 보손으로 남는다.

\begin{figure}[tbp]
\centering
\includegraphics[width=0.95\textwidth]{lecture_144_figure_11.png}
\caption{두 퍼텐셜이 칠판에 다시 등장하며, 멕시코 모자 골짜기에는 반径방향 변위가 표시되어 있다. 안팎으로 움직이는 이 진동이 물리적 힉스 모드다.}
\lectureframe{00:56:30}
\end{figure}

이에 대응하는 다입자 그림도 있다. 진공은 응축이며, 반径방향 진동은 응축의 밀도를 바꾼다. 따라서 힉스 보손은 이 양자 매질 속의 압축파와 비슷하다. 다만 이 파동이 반드시 안정하다고 추론하지 않는다면 유용한 비유다. 보통의 음파도 붕괴할 수 있으며, 힉스 들뜸도 붕괴할 수 있다.

이 차이를 통해 힉스 보손의 발견이 메커니즘의 첫 증거가 아니라 마지막 빠진 조각이었던 이유를 알 수 있다. 대칭이 깨진 진공과 질량이 있는 약한 보손은 이미 간접적으로 드러나 있었다. 밀도에 해당하는 반径방향 파동은 입자로 생성된 뒤 붕괴 생성물을 통해 확인되어야 했다.

\section{붕괴, 생성, 그리고 무거운 페르미온의 이점}

힉스 보손은 페르미온 질량에 비례하는 세기로 페르미온과 결합한다. 힉스 보손은 페르미온과 반페르미온 쌍으로 붕괴할 수 있다.
\begin{equation}
H\longrightarrow f+\bar f,
\end{equation}
같은 도표를 거꾸로 읽으면 생성 과정이 된다.
\begin{equation}
f+\bar f\longrightarrow H.
\end{equation}
따라서 전자와 양전자의 소멸로 힉스 보손을 만드는 것은 원리상 가능하다. 그러나 역사적으로 초기 충돌기는 힉스 보손을 생성할 에너지가 부족했고, 이후의 장치들은 더 근본적인 문제에 부딪혔다. 전자 유카와 결합이 극도로 작다는 점이다. 가벼운 쿼크의 소멸 과정도 마찬가지로 약하다.

\begin{figure}[tbp]
\centering
\includegraphics[width=0.95\textwidth]{lecture_144_figure_12.png}
\caption{칠판에 그린 힉스 붕괴 꼭짓점. 화살표를 거꾸로 읽으면 페르미온과 반페르미온의 융합으로 힉스 보손이 생성된다.}
\lectureframe{00:59:40}
\end{figure}

페르미온이 무거울수록 유카와 결합은 더 강하다. 이 점에서 톱 쿼크가 두드러지는데, 그 질량은 약 $173\,\mathrm{GeV}$, 즉 양성자 질량의 대략 $180$배다. 강의에서는 양성자 질량의 약 $170$배라는 둥근 수치를 사용한다. 그러나 질량이 $125\,\mathrm{GeV}$에 가까운 힉스 보손은 온셸 $t\bar t$ 쌍으로 붕괴하기에는 너무 가볍다.
\begin{equation}
m_H<2m_t.
\end{equation}
또한 유용한 톱 쿼크 빔을 만들 수도 없다. 톱 쿼크는 오래 사는 빔 입자로 강입자화하기 전에 붕괴하기 때문이다.

양자역학은 두 가지 난관을 모두 우회할 방법을 제공한다. 두 글루온은 가상 톱 쿼크 루프를 거쳐 요동한 뒤 힉스 보손으로 나올 수 있다.
\begin{equation}
g+g\longrightarrow H.
\end{equation}
가상 톱 쿼크 선은 실제 톱 쿼크 쌍 생성에 적용되는 온셸 문턱을 만족할 필요가 없다. 양성자에는 글루온이 가득하므로, 대형 강입자 충돌기에서 일어나는 고에너지 양성자-양성자 충돌은 필요한 글루온 쌍을 강렬하게 공급한다.

\begin{figure}[tbp]
\centering
\includegraphics[width=0.95\textwidth]{lecture_144_figure_13.png}
\caption{가상 톱 쿼크 루프를 통한 글루온 융합. 충돌하는 양성자에서 나온 두 글루온이 루프 안에서 만나 힉스 보손을 생성한다.}
\lectureframe{01:06:35}
\end{figure}

이 때문에 양성자 잔해가 복잡한데도 양성자 충돌기가 결정적인 힉스 생산 시설이 되었다. 양성자는 하나의 기본 입자로서 힉스 보손과 결합하는 것이 아니다. 양성자를 이루는 글루온 성분들이 무거운 톱 쿼크의 결합이 지배하는 루프를 통해 힉스 보손을 생성한다.

\section{2012년 입자와 쌍광자 검증}

2012년에 보고된 입자의 질량은 다음 값에 가까웠다.
\begin{equation}
m_H\simeq125\,\mathrm{GeV}.
\end{equation}
강의에서 편의상 비교한 수치로는 양성자 질량의 대략 $127$배다. 이 입자의 발견으로 최소 표준 모형이 요구하는 입자 목록이 완성되었다. 이제 남은 질문은 이 입자의 정량적 성질이 표준 모형이 예측하는 성질과 정확히 일치하는가였다.

톱 쿼크 루프의 논리를 반대 방향으로 적용할 수도 있다. 힉스 보손은 루프 안의 가상 전하 입자들과 결합하고, 그 루프는 광자 두 개를 방출할 수 있다.
\begin{equation}
H\longrightarrow\gamma+\gamma.
\end{equation}
이 붕괴는 드물지만, 광자는 강입자 충돌기에서 배경이 특히 깔끔한 입자다. 따라서 광자의 불변 질량 분포에서 좁은 초과 신호를 찾으면, 더 흔한 붕괴들이 실험적으로 더 복잡한 경우에도 힉스 보손을 찾아낼 수 있다.

\begin{figure}[tbp]
\centering
\includegraphics[width=0.95\textwidth]{lecture_144_figure_14.png}
\caption{칠판 오른쪽의 쌍광자 삼각 루프와 왼쪽의 앞서 살펴본 글루온 융합 루프. 두 과정은 루프 수준의 동일한 논리를 서로 반대 방향으로 살펴본다.}
\lectureframe{01:09:00}
\end{figure}

강의 당시 보고된 쌍광자 붕괴율은 가장 단순한 표준 모형 예측보다 아마도 1.5배 정도 높아 보였으며, 유의도는 약 2 표준편차에 불과했다. 그런 요동만으로는 증거가 되지 않았다. 하지만 그 차이가 계속 관측되었다면, 루프에 추가 전하 입자가 있음을 시사할 수도 있었고, 그 입자는 초대칭이나 표준 모형의 다른 확장과 관련되었을 가능성이 있었다.

이 말은 2012년 당시의 상황에 관한 것이다. 이는 실험물리학자들이 다음에 무엇을 주시해야 한다고 들었는지를 기록한 것이지, 현재의 이상 현상을 주장하는 말이 아니다. 발견 자체는 약한 상호작용 보손 자료가 이미 뒷받침하던 틀을 대체로 확인해 준 결과였다. 더 엄밀한 과학적 검증은 붕괴율에 달려 있었다. 새 스칼라 입자는 힉스 메커니즘이 예측한 질량 의존적 양상에 따라 각 입자와 결합하는가?

\section{청중의 질문}

\begin{classroomqa}
\audiencequestion{왜 어떤 페르미온은 다른 페르미온보다 왼쪽과 오른쪽 사이에서 더 빠르게 진동하며, 따라서 질량도 다른가?\lecturetimestamp{01:12:00}}
\lecturerresponse{페르미온 종류마다 힉스 응축체와의 고유한 결합이 있다. 결합이 강할수록 왼쪽--오른쪽 혼합이 강해지고 질량도 커진다. 같은 결합은 물리적 힉스 보손이 해당 페르미온으로 얼마나 강하게 붕괴하는지도 결정한다. 그래서 일반적으로 무거운 페르미온일수록 더 강하게 결합한다. 표준 모형은 이 유카와 결합들을 매개변수로 나타내지만, 그 뚜렷한 계층 구조가 왜 생기는지는 설명하지 않는다. 전자의 결합은 약하고 톱 쿼크의 결합은 강하다고 말하는 것은 정확한 기술이지만, 그것만으로 더 근본적인 유도가 되는 것은 아니다.}
\end{classroomqa}

\begin{classroomqa}
\audiencequestion{진공에서 힉스 장의 크기를 알고 있나요?\lecturetimestamp{01:13:36}}
\lecturerresponse{그렇다. 여기서 사용하는 규약으로 진공 기대값은 약 $240\,\mathrm{GeV}$이며, 더 정밀한 표준값은 약 $246\,\mathrm{GeV}$이다. 이를 장이 원점에서 벗어난 반지름 방향 변위로 볼 수도 있고, 응축체의 밀도로 볼 수도 있다. 물리적 힉스 양자는 그 값의 작은 반지름 방향 진동이며, 달리 말하면 응축체 밀도의 국소적 요동이다.}
\end{classroomqa}

\begin{classroomqa}
\audiencequestion{힉스 장의 영이 아닌 진공값은 빈 공간의 에너지 밀도를 뜻하는가?\lecturetimestamp{01:14:12}}
\lecturerresponse{그렇다. 힉스 퍼텐셜은 양자장론의 여러 다른 부문과 마찬가지로 진공 에너지에 기여한다. 관측된 순 진공 에너지는 이러한 기여분의 자연스러운 크기에 비해 극히 작으므로, 기여분들은 거의 상쇄되어야 한다. 이 상쇄를 설명하는 이론은 아직 받아들여지지 않았다. 이것이 우주상수 문제이며, 기초 물리학에서 가장 심오한 미해결 문제 가운데 하나다.}
\end{classroomqa}

\section{이 구성이 밝혀낸 것}

이제 끈적한 에테르를 끌어들이지 않고도 논리의 처음부터 끝까지를 살펴볼 수 있다. 양자화된 내부 회전은 전하의 모형을 제공한다. 그릇 모양의 퍼텐셜을 멕시코 모자형으로 바꾸면 장은 영이 아닌 진공값을 갖는다. 불확정성 원리는 고전적으로 정해진 진공을 개별 양자에 민감하지 않은 응축 상태로 바꾼다. 그러면 이 응축 상태가 왼손성 페르미온과 오른손성 페르미온 사이의 약한 상호작용의 양자수 차이를 지닐 수 있어, 두 페르미온이 혼합되고 그 결과 질량을 갖게 된다. 또한 응축 상태는 종방향 모드를 제공하여 $W$와 $Z$가 질량을 가진 벡터 보손으로 전파되게 하며, 광자는 질량이 없는 상태로 남는다.

이러한 재배열 뒤에 남는 반지름 방향 진동이 물리적 힉스 보손이다. 힉스 보손은 ``질량 그 자체''가 아니며, 양성자 질량의 대부분을 만드는 원인도 아니다. 힉스 보손은 대칭 깨짐 장의 관측 가능한 밀도파다. 그 결합에는 해당 장이 만들어 낸 질량의 흔적이 남아 있으며, 그래서 무거운 입자들이 힉스 보손의 생성과 붕괴를 지배한다. 따라서 글루온 융합 루프, 쌍광자 붕괴 경로, 청중이 마지막에 던진 진공 에너지 질문은 서로 별개의 궁금증이 아니다. 이들은 모두 동일한 비자명한 진공이 낳은 결과이자 검증 대상이며, 아직 풀리지 않은 함의다.
